Optical and electrophysiological neural recordings have exponentially increased in scale over the past several decades~\parencite{stevenson2011advances,urai2022large}.
These methods are now maturing into plug-and-play technologies that are commercially available at increasingly affordable rates, raising the possibility of scaling up neural datasets in fundamentally new ways.
In particular, population-level neural recordings can now be obtained across many subjects, brain regions, and repeated behavioral sessions, but summarizing these multi-session datasets is a formidable challenge.
While prior works have introduced relevant methods~\parencite[see][for a review]{klabunde2023similarity}, future advances will require us to further develop and refine these approaches.
Here, we used shape metrics to show that neural population geometry is neither uniform across the brain nor idiosyncratic across animals, but structured at both scales: organized by anatomy across regions and by behavior and genotype across individuals.

% Compared with existing methods, we highlight two advances.
% First, we propose notions of distance that are satisfy the triangle inequality and are symmetric (unlike, e.g., linear predictivity scores~\hl{fig}).
% Such properties are desirable for many downstream analysis methods.
% For example, the notion of a ``cluster'' of datapoints is ambiguous when the triangle inequality is badly violated.
% Second, our approach leverages explicit geometric transformations (e.g. rotations and reflections) to quantify similarity.
% This forces practitioners to be explicit with their modeling assumptions and enables a flexible framework for model comparison and alignment.

Systematic comparisons of neural population recordings will require a variety of approaches including both unsupervised methods for exploratory analysis (e.g. dimensionality reduction and clustering) as well as supervised methods for testing targeted hypotheses (e.g. regression and classification models).
We showed that a number of these analyses are enabled by establishing notions of distance that are symmetric and satisfy the triangle inequality (i.e. define a metric space).
For example, we used clustering in Procrustes shape space to identify brain regions with similar neural population activity, and found that these clusters were in line with anatomical knowledge (Fig. \ref{fig:allen},\ref{fig:miller}).
It is now widely appreciated that many task and behavioral variables can be decoded from nearly any recorded brain region [\hl{cite: Stringer, Musall, Steinmetz, or the IBL brainwide map}]. This has been read as evidence for a highly distributed code, but decodability of a single variable is a coarse summary of a population's activity. Our results suggest a complementary view: when regions are compared at the level of population geometry rather than decodability, a clear anatomical organization re-emerges, with shape distance recovering a hierarchy across dozens of regions without any anatomical labels supplied to the algorithm (\textit{Fig.} \ref{fig:allen}, \ref{fig:miller}).

We also demonstrated that shape metrics can be leveraged to fit supervised models.
For example, a neural system's position within Procrustes shape space can be used to predict single-neuron properties (Fig. \ref{fig:hdcells}), anatomical location (Fig. \ref{fig:miller}) and  psychometric performance (Fig. \ref{fig:ibl}).
Most strikingly, individual differences in behavior were mirrored by individual differences in population geometry extending to full cohorts a link that prior work had established only for pairs of subjects.
Recent work by \textcite{fascianelli2024neural}, which performed a detailed analysis of two nonhuman primate subjects and found correlations between their behavioral differences and neural geometries.
Shape metrics provide a rigorous framework to scale up these investigations to larger experimental cohorts (e.g. across \hl{20 mice} in Fig. \ref{fig:ibl}). Crucially, this framework is agnostic to recording modality. For instance, applied to human EEG, shape distances reveal reliable subject-specific geometry, and show that self-supervised pretraining of a foundation model \parencite{cui2024,wang2025} amplifies rather than averages out these idiosyncrasies \parencite{marty2026self}.

Shape metrics are not the only family of distance functions that satisfy symmetry and the triangle inequality.
Slight modifications to CKA scores and CCA-based similarity scores result in proper metric spaces~\parencite{Williams2021}.
Distance-based RSA scores can also be turned in proper metrics, as can be seen from recent work showing equivalence between RSA and CKA~\parencite{Williams2024}. These alternatives are complementary to ours: the contribution here is not the existence of a metric per se, but the demonstration of what metric structure buys on neurobiological data: clustering, regression, and nearest-neighbor prediction across large collections of recordings.

Linear predictivity scores are another framework for comparing neural systems.
They are particularly popular for quantifying similarity between artificial and biological networks \parencite[e.g.][]{conwell2022can}.
One can view shape metrics as a framework for incorporating constraints into predictivity scores in order to turn them into proper metric spaces.
Indeed, raw predictivity scores can be arbitrarily asymmetric---it is easy to construct a synthetic example such that predicting neural system $A$ from $B$ yields $R^2 \approx 1$ while predicting neural system $B$ from $A$ yields $R^2 \approx 0$ (see simulations in \textit{Fig.} \ref{fig:synth}f). Since there is often no \textit{a priori} asymmetry between the systems we would like to collectively evaluate, it is natural to compare them within a proper metric space.
However, asymmetry is not strictly a defect.
In fact, it is desirable in circumstances where one only cares about predicting $A$ from $B$ (and not $B$ from $A$).
This might be the case, for example, when artificial deep networks are used to predict biological responses for the purpose of running \textit{in silico} experiments. Decoding analyses have a related advantage. Because a decoder targets a specific variable, it answers a question that shape distances cannot: whether a particular piece of information is present in a population, and in a format that a downstream reader could exploit. A shape distance instead summarizes the geometry as a whole, which is the appropriate summary when the question concerns the structure of a representation rather than the accessibility of one variable within it. The two are naturally used together and indeed we relied on decoding and regression throughout this work to validate and interpret the structure that shape distances revealed.


% While the flexibility of the shape metrics framework can be advantageous, it also highlights a challenge.
% Shape metrics represent only a fraction of a broader array of methodologies---a recent review by \textcite{klabunde2023similarity} cataloged over thirty different representational similarity measures from the machine learning literature.
% Many of these could, in principle, be applied to neurobiological data.
% How should practitioners choose among the large collection of methods that quantify similarity in neural population codes?

% By defining similarity in terms of explicit geometric alignment transformations (e.g. rotations or permutations of neural firing rate space), this framework forces practitioners to be explicit with their modeling assumptions.
% Furthermore, it enables us to reason about and compare different analysis outcomes.
% For example, we know that the Procrustes distance between two systems will \textit{always} be less than or equal to the one-to-one matching distance because the former allows for a more permissive set of alignment transforms.
% Neuroscientists can use these relationships among shape metrics to interpret and compare different forms of analysis.

For simplicity, we utilized the Procrustes distance as a running example throughout most of our analyses.
However, this is just one of a broader class of possible shape metrics.
Our practical experience suggests that the Procrustes distance is reasonable default choice---it is resistant to overfitting and has no hyperparameters that need to be tuned---but practitioners should consider a broader set of distance functions depending on their analysis goals.
For example, the Procrustes distance is not sensitive to differences in the tuning structure of individual neurons, as we demonstrated in Fig. \ref{fig:synth}e-g.
A more stringent distance measure is given by the one-to-one matching distance, \cref{eq:one-to-one-matching}, and its generalization, the \textit{soft matching distance}~\parencite{khosla2024soft}.
These alternatives can be used to quantify the extent to which neural systems represent information in arbitrarily rotated coordinate systems or in terms of a privileged neural basis %reproducible basis vectors 
\parencite{Khosla2024}.

By grounding comparisons of neural systems in explicit alignment transformations, shape metrics are related to a number of previous analyses.
For example, \textcite{Chen2021} analyzed hippocampal population activity and found that they could use rotational alignments to predict neural responses in one condition (e.g. right turns on a T-maze) using cross-subject alignments on a different condition (e.g. left turns).
Similar cross-condition decoding analyses have been applied elsewhere \parencite{saez, bernardi, barbosa2023}.
Here we showed that such rotational alignments can be used to define a metric space, and therefore used to rigorously analyze larger collections of neural systems (not just pairs of subjects or conditions).
In contrast, it is more difficult to draw direct connections to alternative representational similarity measures, like CKA and RSA, which are not defined in terms of explicit geometric transformations \parencite[but see][for a relation between Procrustes distance and RSA-style methods]{harvey2024duality}.

A further consequence of defining similarity through explicit alignment transformations is that the alignments themselves are useful objects. Once a collection of neural systems has been aligned to a common reference (by generalized Procrustes analysis, in our case) one obtains a shared coordinate system in which all subjects can be jointly analyzed. We exploited this to decode head direction in a held-out mouse from a decoder fit entirely on other mice (\textit{Fig.} \ref{fig:hdcells}e), and to predict a subject's psychometric curve from its nearest neighbors in shape space (\textit{Fig.} \ref{fig:ibl}f). This links our framework directly to hyperalignment \parencite{haxby2020hyperalignment}, which pursues the same goal in human neuroimaging, and to latent-space alignment methods developed for across-session and across-subject decoding [\hl{cite: Degenhart, Gallego, or your preferred latent-alignment refs}]. Distance and alignment are two aspects of the same computation: the distance quantifies how much structure two systems share, while the alignment supplies the map that makes shared structure usable. Similarity measures that return only a scalar score do not furnish this map.

The potential of the shape metrics framework to study neural geometry was not fully explored here.
We have focused on analyzing trial-averaged neural responses, but the scale and orientation of trial-to-trial noise is important to theories of neural coding \parencite{moreno2014information}.
Furthermore, there is substantial interest in moving beyond geometry to understand neural circuits as dynamical systems \parencite{vyas2020computation}.
The framework we've outlined can be extended to address these issues \parencite{duong2023representational,ostrow2024beyond,nejatbakhsh2024comparing}.
However, these extensions remain a relatively active area of research.

We focused primarily on Procrustes shape distance, which has no tunable hyperparameters and is both analytically and computationally tractable.
Leveraging the metric space properties of this distance, we revealed several novel and non-trivial insights. Overall, our results reveal that neural population codes carries structure at multiple scales: across-subject differences in head direction codes, an anatomical hierarchy across regions recovered without supervision, and a correspondence between population geometry and behavior
As recording technologies continue to advance, generating increasingly large and complex neural datasets, shape metrics have the potential to become a fundamental tool for neuroscience, enabling researchers to systematically analyze large collections of neural systems at the level of population activity.


% \hl{something about how explicit alignments can be used for across-animal decoding}



% Shape metrics are a family of distance functions that leverage explicit geometric transformations---rotations and reflections in the case of Procrustes distance, and permutations in the case of one-to-one matching distance---to quantify similarity in neural representations.
% % This forces practitioners to be explicit with their modeling assumptions and enables a flexible framework for model comparison and alignment.

% By placing multiple measures of distance into a common framework, we enable direct comparisons between these analyses.
% For example, the Procrustes distance between two systems is a lower bound on the soft matching distance between them, since the latter incorporates additional constraints into the alignment transform.


% Shape metrics are based on explicit alignment transformations

% \hl{clean up}
% The second advantage of shape metrics is  the use of explicit geometric transformations to align and compare neural systems.
% In order test different scientific hypotheses, the transform class can be modified.
% For example, the Procrustes distance between two systems can be thought of as the root-mean-squared error obtained after an optimal translation, rotation, and reflection to align the firing rate spaces.
% Intuitively, this quantifies similarity in manifold shape in a way that does not depend on the choice of coordinate axes.

% A more stringent notion of similarity is given by the permutation Procrustes distance and its generalization, \textit{soft matching distance}~\parencite{khosla2024soft}.
% This concept was used in a recent study to demonstrate that certain neural representations are not randomly oriented---instead, similar tuning statistics at the single unit level are found across neural systems over a broad range of tasks \parencite{Khosla2024}.
% By placing multiple measures of distance into a common framework, we enable direct comparisons between these analyses.
% For example, the Procrustes distance between two systems is a lower bound on the soft matching distance between them, since the latter incorporates additional constraints into the alignment transform.


% The potential of the shape metrics framework to study neural geometry was not fully explored here. For example,  recent work has leveraged the so-called cross-condition decoding to compare the encoding geometry across different conditions \parencite{saez, bernardi, barbosa2023}. One specific application of this approach is to find abstract encoding in the brain --- i.e. encoding of one condition that is invariant to other conditions. Decoding performance, like regression, is asymmetric and it saturates at chance and maximum accuracies. Indeed, decoding often leads to asymmetric accuracies which are hard to interpret in the context of abstract representations (e.g. some condition C is considered more abstract when trained along with condition A and tested along with condition B, than vice-versa). One approach is to average out these asymmetries across all conditions \parencite{bernardi}, but future work could consider using shape metrics to find abstract representations. Other applications have been explored elsewhere. For example, while our study focused on deterministic shape metrics, stochastic extensions provide a more nuanced understanding of variability in neural representations \parencite{duong2023representational}. \hl{. what else? non linear methods?}

% Overall, shape metrics provide a robust and versatile framework for quantifying differences in neural population activity, offering new insights into the complex interplay between representational geometry and behavior. As recording technologies continue to advance, generating increasingly large and complex neural datasets, shape metrics have the potential to become a fundamental tool for neuroscience, enabling researchers to systematically compare brain representational architectures across animals, potentially even of species.


% % We have showcased shape metrics here by applying it to several neural datasets collected from rodents and non-human primates, demonstrating its potential to measure individual variability \hl{(cite Fig)}, reveal across-region representational similarities \hl{(cite Fig)} as well as how differences in neural geometry captured by shape analyses are reflected in different behavioral strategies \hl{(cite Fig)}. With these analyses, we highlight the advantage of rigorously calculate distances between neural populations. For instance, in our analysis of head direction cells in the mouse postsubiculum, shape metrics effectively captured the consistent ring-like structure encoding head direction angles, but we also show how shape metrics can be leveraged to reveal individual differences. Through conservative cross-validation methods, we show that individual differences can be captured by shape metrics and are reflected in the behavioral strategies. 